\section*{अध्याय \ref{ch:b2:chromatography} --- वर्णलेखन}

\begin{solution}{exo:b2:chromatography:1}
$k = (5.0 - 1.0)/1.0 = 4.0$; \omterm{def:b2:chromatography:principle}{गतिशील प्रावस्था} में समय का अंश $1/(1 + k) = 0.20$।
\end{solution}

\begin{solution}{exo:b2:chromatography:2}
$N = 16(6.0/0.30)^2 = 16 \times 400 = 6400$।
\end{solution}

\begin{solution}{exo:b2:chromatography:3}
$H = \qty{250}{mm}/10\,000 = \qty{0.025}{mm} = \qty{25}{\micro m}$।
\end{solution}

\begin{solution}{exo:b2:chromatography:4}
रक्त में एथेनॉल: GC (वाष्पशील)। प्रोटीन: HPLC। पेय में कैफ़ीन: HPLC (जल में, बिना तैयारी के पर्याप्त
वाष्पशील नहीं)। पेट्रोल में बेन्ज़ीन: GC। शर्कराएँ: HPLC (वे क्वथन से पहले विघटित हो जाती हैं)।
\end{solution}

\begin{solution}{exo:b2:chromatography:5}
$R_s = 2(4.6 - 4.0)/(0.40 + 0.50) = 1.2/0.90 \approx 1.3$: पूरी तरह आधार-रेखा तक पृथक नहीं (1.5 से कम)।
\end{solution}

\begin{solution}{exo:b2:chromatography:6}
$\sqrt N = 4R_s\,\dfrac{\alpha}{\alpha - 1}\,\dfrac{1 + k_2}{k_2} = 4 \times 1.5 \times 11 \times 1.25 =
82.5$, अतः $N \approx 6.8 \times 10^3$।
\end{solution}

\begin{solution}{exo:b2:chromatography:7}
$u_{\mathrm{opt}} = \sqrt{8/2} = \qty{2}{mm/s}$; $H_{\min} = 5 + 2\sqrt{16} = \qty{13}{\micro m}$।
\end{solution}

\begin{solution}{exo:b2:chromatography:8}
पहले फ़ीनॉल (ध्रुवी \ce{OH}, जल के साथ हाइड्रोजन आबंध), फिर बेन्ज़ीन, फिर टॉलूईन (एक अतिरिक्त
\ce{CH3}, अधिक अध्रुवी)। अधिक मेथेनॉल \omterm{def:b2:chromatography:principle}{गतिशील प्रावस्था} को कम ध्रुवी बनाता है: तीनों पहले उत्क्षालित होते हैं,
उसी क्रम में।
\end{solution}

\begin{solution}{exo:b2:chromatography:9}
$F = (860/400)/(20.0/10.0) = 2.15/2.00 = 1.075$। नमूना: $c_X = (645/410) \times 10.0/1.075 \approx
\qty{14.6}{mg/L}$।
\end{solution}

\begin{solution}{exo:b2:chromatography:10}
$R_s \propto \sqrt N \propto \sqrt L$: $L$ को $(1.5/1.06)^2 \approx 2.0$ से गुणा करना होगा, जो विश्लेषण समय
और दाब को दुगुना कर देता है। अन्य उत्तोलक: वरणात्मकता (\omterm{def:b2:chromatography:principle}{गतिशील प्रावस्था}, \omterm{def:b2:chromatography:principle}{स्थिर प्रावस्था},
ताप), छोटे कण, इष्टतम प्रवाह दर।
\end{solution}

\begin{solution}{exo:b2:chromatography:11}
शिखर $4\sigma R_s$ की दूरी पर हैं, अतः मध्य-बिंदु हर एक से $2\sigma R_s$ पर है। $R_s = 1$: हर शिखर
$\mathrm e^{-(2)^2/2} = \mathrm e^{-2} = 0.135$ का योगदान देता है; घाटी शिखर ऊँचाई के $0.27$ पर है।
$R_s = 1.5$: $2\mathrm e^{-(3)^2/2} = 2\mathrm e^{-4.5} \approx 0.022$, लगभग 2~\%: व्यावहारिक रूप से
आधार-रेखा पर।
\end{solution}

\begin{solution}{exo:b2:chromatography:12}
$k_2/(1+k_2)$: 0.67, 0.83, 0.91; $t_M$ की इकाइयों में समय: 3, 6, 11। 2 से 5 पर जाने से दुगुने समय के बदले
\omterm{def:b2:chromatography:separation}{विभेदन} में 25~\% लाभ होता है; 5 से 10 पर, लगभग फिर दुगुने समय के बदले 9~\%। लगभग 2 से 5 तक का $k$
सामान्य समझौता है।
\end{solution}

\begin{solution}{pb:b2:chromatography:1}
\textbf{1.} स्थिर: लंबी ऐल्किल शृंखलाएँ धारण करने वाले सिलिका कण, अध्रुवी। गतिशील: मेथेनॉल के साथ
जल, ध्रुवी।
\textbf{2.} वे अत्यंत ध्रुवी हैं और \omterm{def:b2:chromatography:principle}{गतिशील प्रावस्था} में रहती हैं: $k \approx 0$।
\textbf{3.} कैफ़ीन में थियोफ़िलीन से एक अधिक मेथिल समूह है, और थियोफ़िलीन में एक \ce{N-H} है जो
जल के साथ हाइड्रोजन आबंध बनाता है: थियोफ़िलीन अधिक ध्रुवी है और पहले उत्क्षालित होती है।
\textbf{4.} यह संरचना में निकट है (समान अनुक्रिया और व्यवहार), पेय में अनुपस्थित है, और कैफ़ीन के
निकट उत्क्षालित होते हुए भी उससे विभेदित है।
\textbf{5.} दोनों घटेंगे: कम ध्रुवी \omterm{def:b2:chromatography:principle}{गतिशील प्रावस्था} \omterm{def:b2:chromatography:principle}{स्थिर प्रावस्था} से बेहतर प्रतिस्पर्धा करती है।
\textbf{6.} \omterm{def:b2:chromatography:principle}{गतिशील प्रावस्था} को स्तंभ पार करने में लगने वाला समय, अर्थात् वह समय जो कोई यौगिक चलते हुए बिताता है।
\textbf{7.} $k_{\mathrm{theo}} = (3.0 - 1.0)/1.0 = 2.0$; $k_{\mathrm{caf}} = 2.4$।
\textbf{8.} $\alpha = 2.4/2.0 = 1.2$।
\textbf{9.} कैफ़ीन $N = 16(3.4/0.20)^2 = 4624 \approx 4.6 \times 10^3$; थियोफ़िलीन $16(3.0/0.18)^2
\approx 4.4 \times 10^3$।
\textbf{10.} $H = \qty{150}{mm}/4624 \approx \qty{0.032}{mm} = \qty{32}{\micro m}$।
\textbf{11.} $R_s = 2(3.4 - 3.0)/(0.18 + 0.20) = 0.80/0.38 \approx 2.1$।
\textbf{12.} $\frac{\sqrt{4624}}{4} \times \frac{0.2}{1.2} \times \frac{2.4}{3.4} = 17 \times 0.167 \times
0.706 \approx 2.0$; छोटा अंतर समीकरण की समान-चौड़ाई कल्पना से आता है।
\textbf{13.} हाँ: $R_s > 1.5$।
\textbf{14.} $N$ आधा हो जाता है: $R_s \approx 2.1/\sqrt2 \approx 1.5$, ठीक आधार-रेखा पृथक्करण पर।
\textbf{15.} आधा: कैफ़ीन \qty{1.7}{min} पर।
\textbf{16.} $R_s \approx 2.1\sqrt2 \approx 3.0$, दुगुने समय और दुगुने दाब पर: यहाँ यह व्यर्थ है।
\textbf{17.} \omterm{def:b2:chromatography:principle}{गतिशील प्रावस्था} का संघटन (मेथेनॉल अंश, pH, कोई अन्य कार्बनिक विलायक) या
\omterm{def:b2:chromatography:principle}{स्थिर प्रावस्था}।
\textbf{18.} थोड़ी: $k_2/(1 + k_2) = 0.71$ पहले से है; $k_2 = 5$ पर यह 0.83 (+18~\%) होता, जबकि
विश्लेषण समय $6/3.4 \approx 1.8$ से गुणा हो जाता।
\textbf{19.} $F = (1500/1000)/(50.0/40.0) = 1.50/1.25 = 1.20$।
\textbf{20.} दोनों शिखर एक ही अंतःक्षेपण से आते हैं: क्षेत्रफलों के अनुपात में आयतन कट जाता है।
\textbf{21.} $c = (970/1010) \times 40.0/1.20 \approx \qty{32.0}{mg/L}$।
\textbf{22.} दस गुना तनुकरण: \qty{320}{mg/L}।
\textbf{23.} $A_{IS}$ बहुत बड़ा होता, और कैफ़ीन का परिणाम बहुत कम।
\textbf{24.} नमूने जैसी ही परिस्थितियों में अंतःक्षेपित कैफ़ीन मानकों की एक श्रृंखला,
सांद्रता के विरुद्ध क्षेत्रफल की अंशांकन रेखा, और पुनरुत्पादनीय अंतःक्षेपण आयतन।
\textbf{25.} $\qty{320}{mg/L} \times \qty{0.250}{L}$: प्रति डिब्बा कैफ़ीन का $\boldsymbol{\approx \qty{80}{mg}}$ (अभ्यास के
आँकड़े)।
\end{solution}
